\section{Formal semantics}\label{sec:formal}

\newcommand{\keval}{\Rightarrow }
\newcommand{\qeval}{\mapsto_q }
\newcommand{\ceval}{\mapsto_c }

Before presenting a formal operational semantics for qstack programs, we briefly review the relevant standard mathematical concepts and notation of quantum computation that we use. This section serves as a reference to make the paper self-contained for readers from programming languages backgrounds who may be less familiar with quantum computing formalisms. The concepts presented below are well-established in the quantum computing literature.

\subsection{Quantum Programs Basic Concepts}

A qubit is a unit vector in the 2-dimensional complex Hilbert space, $\mathbb{C}^2$. Its general form is:
\[
\ket{\psi} = \alpha\ket{0} + \beta\ket{1}, \quad \alpha, \beta \in \mathbb{C}, \quad |\alpha|^2 + |\beta|^2 = 1
\]

Unlike classical bits with definite values of 0 or 1, qubits encode information through quantum superposition—a linear combination of basis states $\ket{0}$ and $\ket{1}$. The complex coefficients $\alpha$ and $\beta$ determine not only the probabilities $|\alpha|^2$ and $|\beta|^2$ of measuring the corresponding basis states, but also contain phase information that affects quantum interference patterns during computation. This phase information, while unobservable in direct measurements, plays a crucial role in quantum algorithms through constructive and destructive interference.

Quantum operations can affect both the probability amplitudes and phases. For example, a phase flip operation transforms a state $\ket{+} = \frac{1}{\sqrt{2}}(\ket{0} + \ket{1})$ into $\ket{-} = \frac{1}{\sqrt{2}}(\ket{0} - \ket{1})$ without altering the measurement probabilities, yet fundamentally changing the quantum state's behavior in subsequent operations.

An $n$-qubit system is described by a state $\ket{\psi^n} \in \mathbb{C}^{2^n}$, which lives in the tensor product of $n$ single-qubit Hilbert spaces. When the dimension of the space is not relevant, we denote the state simply as $\ket{\psi}$. Adding a new qubit in state $\ket{0}$ at position $n$ results in:
\[
\ket{\psi^{n+1}} = \ket{\psi^n} \otimes \ket{0}
\]

The computational basis of an $n$-qubit system is the set of $2^n$ orthonormal states $\{\ket{x} : x \in \{0,1\}^n\}$, where each $\ket{x}$ corresponds to a classical bitstring $x = (x_0, x_1, \ldots, x_{n-1})$. These basis states represent configurations in which each qubit is either in state $\ket{0}$ or $\ket{1}$, and they span the space $\mathbb{C}^{2^n}$. Any quantum state $\ket{\psi^n}$ can be expressed as a linear combination of these:

\[
\ket{\psi^n} = \sum_{x \in \{0,1\}^n} \alpha_x \ket{x}
\]


A unitary operator $U$ is a linear transformation such that $U^\dagger U = UU^\dagger = I$. Applying a unitary to a quantum state preserves its norm and produces another valid quantum state. The composition of two unitaries is also unitary.

To apply a unitary $U$ that acts on a single qubit $j$ in an $n$-qubit system, we define a lifted operator $U^{(j)}$ that acts as the identity on all qubits except at position $j$:
\[
U^{(j)} = I \otimes \cdots \otimes I \otimes U \otimes I \otimes \cdots \otimes I
\]
with $U$ appearing in the $j$-th position.

If $U$ acts on $k$ qubits, we assume it is defined to act on the last $k$ positions of a tensor product space. To apply it to any arbitrary set of $k$ qubits at positions $j_1, j_2, \ldots, j_k$, we use a permutation operator $\Pi$ that moves those qubits to the last $k$ positions. The lifted operator is then defined as:
\[
U^{(j_1, \ldots, j_k)} = \Pi^{-1} \circ (I \otimes \cdots \otimes I \otimes U) \circ \Pi
\]
The permutation operator $\Pi$ is unitary, since it rearranges tensor product factors without changing inner products. Therefore, $U^{(j_1, \ldots, j_k)}$ is also unitary.

The state evolves under the $U$ operation as:
\[
\ket{\psi'} = U^{(j_1, \ldots, j_k)} \ket{\psi^n}
\]


Measurement of qubit $j$ in the computational basis is denoted by $\|\psi^n\|_j \mapsto b, \ket{\psi^{n-1}}$. This operation produces a classical outcome $b \in \{0, 1\}$ and a post-measurement state $\ket{\psi^{n-1}}$, calculated as follows: assume $\ket{\psi^n}$ is expressed in the computational basis as:
\[
\ket{\psi^n} = \sum_{x \in \{0,1\}^n} \alpha_x \ket{x}
\]
where $x = (x_0, x_1, \ldots, x_{n-1})$ is a bitstring of length $n$. Then the probability of observing outcome $b$ is:
\[
P(b) = \sum_{\substack{x \in \{0,1\}^n \\ x_j = b}} |\alpha_x|^2
\]
and the post-measurement state is:
\[
\ket{\psi^{n-1}} = \frac{1}{\sqrt{P(b)}} \sum_{\substack{x \in \{0,1\}^n \\ x_j = b}} \alpha_x \ket{x}
\]
This state lies in the subspace where qubit $j$ is fixed to $b$, reducing the dimensionality of the space from $2^n$ to $2^{n-1}$. Intuitively, we think of measurement as also deallocating the measured qubit: after measuring, we no longer treat it as part of the quantum state.

\subsection{Formal Semantics of qstack Programs}

To formalize the behavior of our kernel-based quantum language, we define a small-step operational semantics that captures the interplay between classical control and quantum state evolution. Execution of a kernel introduces a new qubit, performs a sequence of operations, and concludes with a callback function that governs the next step based on measurement outcomes. The model defines a stack discipline for qubit allocation and use, supporting dynamic control flow through measurement-driven callbacks.

The formal syntax of qstack programs is given by the grammar below:


{\tt\small
\begin{tabular}{lrl}
program                 &:==& kernel \\
kernel                  &:==& (\textit{qubit\_id}:instructions:callback) \\
instructions            &:==& instruction, instructions | $\varnothing$ \\
instruction             &:==& quantum\_gate | kernel \\
quantum\_gate           &:==& \textit{gate\_id}(targets) \\
targets                 &:==& \textit{qubit\_id} | \textit{qubit\_id}, target  \\
callback                &:==& \textit{callback\_id} \\
kernel\_option          &:==& kernel | $\varnothing$ \\
\end{tabular}}

Each gate identifier \textit{gate\_id} is assumed to denote a fixed unitary operator of appropriate arity. We write $U_{\mathtt gate\_id}$ to refer to the unitary associated with a particular gate identifier. Similarly, each callback identifier \textit{callback\_id} denotes a classical function of type $\mathtt{bool} \times \mathtt{State} \rightarrow \texttt{kernel\_option} \times \mathtt{State}$, which we write as $f_\mathtt{callback\_id}$. We assume this association is fixed and globally defined. We use $S$ to range over \texttt{State} which is simply arbitrary classical state we formally leave opaque.

For simplicity, this formal semantics does not explicitly encode the notion of instruction sets or compilation layers. While qstack programs in practice operate within specific instruction sets (groups of related quantum operations) and can be compiled between different instruction sets or enhanced with error correction, the core operational behavior can be understood through this simplified model where gates are globally available and are treated as abstract unitary operators. 

A program consists of a single kernel. We formally represent each kernel as a triple $(q : I : cb)$, where:
\begin{itemize}
    \item $q$ is a fresh logical qubit identifier.
    \item $I$ is a finite sequence of instructions to execute on that qubit and others in scope.
    \item $cb$ is a callback id that takes a measurement outcome and classical state, and produces a new kernel.
\end{itemize}

An instruction is either a quantum gate or a nested kernel.

Execution is modeled via a small-step relation on configurations $(C, S, (\ket{\psi}, m))$, where:
\begin{itemize}
    \item $C$ is a stack of instruction sequences and associated callbacks (essentially the call stack of kernel invocations).
    \item $S$ is the classical store (e.g., variables and memory).
    \item $(\ket{\psi}, m)$ is the current quantum state vector $\psi$ and a mapping $m$ from logical qubit names to positions in the state vector.
\end{itemize}

We denote the evaluation with:
$$
C, S, (\ket{\psi}, m) \keval C', S', (\ket{\psi'}, m')
$$

Given a program $(q : I : cb)$, we start in the initial configuration $((q : I : cb):[]):[], S_0, ([], [])$ for some $S_0$. The first step applies the \textsc{start-kernel} rule to allocate the initial qubit and begin execution of the top-level kernel. A final state has the form $[]; S; ([], [])$ and we take the outcome of the program to be the final classical state $S$.

The operational semantics is defined by the following inference rules. Which rule applies at each step depends on the form of the instruction sequence at the top of the stack:
\begin{itemize}
\item If the next instruction is a nested kernel, the \textsc{start-kernel} rule applies.
\item If it is a gate application, the \textsc{quantum-gate} rule applies.
\item If there are no instructions left, one of the two \textsc{end-of-kernel} rules applies depending on the callback result.
\end{itemize}

\subsection*{Kernel Initiation (\textsc{start-kernel})}

\begin{center}
\begin{mathpar}
\inferrule[start-kernel]
{  }
{ ((q_k:I_k:cb_k): I, cb):C', S, (\ket{\psi^n}, m)  \keval (I_k, cb_k):(I, cb):C', S, (\ket{\psi^{n}} \otimes \ket{0},\ m[q \mapsto n])  }

\end{mathpar}
\end{center}

This rule initiates the execution of a kernel by allocating a new qubit in the $|0\rangle$ state and extending the qubit mapping.

A fresh qubit is allocated at the next available position $n$ in the state vector. The logical identifier $q_k$ is mapped to this index, and the quantum state is extended with $|0\rangle$. The instructions $I_k$ are pushed onto the call stack along with their callback $cb_k$.

\subsection*{Quantum Gate Application (\textsc{quantum-gate})}


\begin{center}
\begin{mathpar}
\inferrule[quantum-gate]
{  }
{ (\mathtt{gate\_id}(q_0, \ldots, q_j): I', cb):C', S, (\ket{\psi}, m) \keval (I', cb):C', S, (U_\mathtt{gate\_id}^{(m[q_0], \ldots, m[q_j])}\ket{\psi}, m)  }
\end{mathpar}
\end{center}

This rule handles the application of a quantum gate to one or more target qubits.

If the next instruction is a gate application $\mathtt{gate\_id}(q_0, \ldots, q_j)$, we resolve each $q_i$ to its corresponding index in the quantum state using the mapping $m$, then apply the appropriate lifted unitary operator $U_\mathtt{gate\_id}$ to the current state vector. The call stack is updated to reflect the remaining instructions, and the classical state and qubit mapping are unchanged.

\subsection*{Kernel Completion with Termination (\textsc{end-of-kernel-done})}

\begin{center}
\begin{mathpar}
    
\inferrule[end-of-kernel-done]
{  
    \|\psi^{n}\|_{n-1} \mapsto b, \ket{\psi'} \\ 
    m =   m' \uplus  [q \mapsto n-1]  \\
    f_{cb}(b,S) \mapsto \varnothing, S'  }
{ ([], cb):C', S, (\ket{\psi^n}, m) \keval C', S', (\ket{\psi'}, m') }

\end{mathpar}
\end{center}

This rule completes a kernel when there are no more instructions to execute, and the callback returns no further computation. The last qubit (at position $n - 1$) is measured, yielding a classical bit $b$ and a collapsed state $\ket{\psi'}$. The mapping $m$ is updated to remove the deallocated qubit: $m = m' \uplus [q \mapsto n - 1]$.

\subsection*{Kernel Completion with Continuation (\textsc{end-of-kernel-continue})}


\begin{center}
\begin{mathpar}

\inferrule[end-of-kernel-continue]
{ \|\psi^{n}\|_{n-1} \mapsto b, \ket{\psi'} \\
    m =   m' \uplus  [q \mapsto n-1]  \\
    f_{cb}(b,S) \mapsto (q_k:I_k:cb_k), S'  }
{ ([], cb):C', S, (\ket{\psi^n}, m) \keval (I_k, cb_k):C', S',  (\ket{\psi'} \otimes \ket{0},\ m'[q_k \mapsto n-1])  }

\end{mathpar}
\end{center}

This rule completes a kernel when the callback returns a new kernel to execute next. It is essentially the composition of completing a kernel and starting a different kernel.
